\section{Localization with triangular recenterings}
\label{loc:section}

An open kernel alone does not control the component reached by a bounded
path.  The additional structure used here is a family of real
recenterings, available at every bounded real displacement.  We prove the
required upper containment, including the uniform path selection needed
for the analytic-disc argument.

\label{loc:kernel}
For open sets $D_t\subset\C^d$, $0<t<t_0$, write $\mathcal K(D_t)$
for their open kernel in the sense of Definition~\ref{prep:kernel}.
If a fixed point $p$ belongs to this kernel, let $P$ be its connected
component.  For $L>|p|$, let $A_t(L)$ be the component containing
$p$ in $D_t\cap B(0,L)$, whenever $p\in D_t$.  Thus $P$ is defined
by eventual neighborhood inclusion, whereas $A_t(L)$ records access
from the seed at one fixed scale and radius.

Throughout this section, subanalytic families are restricted to bounded
coordinate boxes and are globally subanalytic there.  In particular,
quantification, closure, connected components in families, and definable
selection are available.  A generic conclusion permits deleting a
countable union of subanalytic sets of dimension smaller than that of the
real parameter box.  A separate exceptional set is allowed for each
integer $L$ and each positive rational accuracy.

\subsection{Normalized domains and recenterings}

Fix a decomposition $\mathbb R^d=\mathbb R^{d_1}\times\cdots
\times\mathbb R^{d_s}$, and write $u_{>j}=(u_{j+1},\ldots,u_s)$.
Let $U\subset\mathbb R^d$ be a real box in a complex coordinate box
$\Omega$.  Suppose $S\subset\Omega$ is a fixed subanalytic open set,
locally pseudoconvex at its boundary in $\Omega$.  In particular $S$
does not change when its real center changes.  Let
$R_t(x)\in\operatorname{GL}_d(\mathbb R)$ be subanalytic and tend to zero,
locally uniformly at generic $x$.  Define
\begin{equation}\label{loc:normalized}
 D_t^x=\{Z:x+R_t(x)Z\in S\},\qquad
 C_t(x,u)=R_t(x)^{-1}R_t(x+R_t(x)u).
\end{equation}
Only portions whose original coordinates lie in $\Omega$ are used.
Assume that, locally at generic $x$, for every finite $M$ and every
integer $k\geq0$,
\begin{equation}\label{loc:ratio}
 \|C_t(x,\cdot)-G_x\|_{C^k(\{|u|\leq M\})}\longrightarrow0,
 \qquad
 G_x(u)=\operatorname{diag}
       (G_{x,1}(u_{>1}),\ldots,G_{x,s}(u_{>s})),
\end{equation}
where all entries of $G_x$ are real polynomials,
$G_x(0)=I$, and $G_x(u)$ is invertible for every real $u$.
The last condition means, in particular, that $G_x^{-1}$ is bounded
on each of the fixed finite boxes that will be used.  Suppose a fixed
$p\in\mathbb C^d$ belongs to $\mathcal K(D_t^x)$ throughout the real
parameter box under consideration.  Put $P_x$ for its marked component.

\begin{theorem}[Accessible upper containment]\label{loc:theorem}
Under these hypotheses, at generic $x$, for every finite $L>|p|$ and
every $\eta>0$ there is $t_0=t_0(x,L,\eta)>0$ such that
\begin{equation}\label{loc:upper}
 0<t<t_0\quad\Longrightarrow\quad
 A_t^x(L)\subset\{Z:\operatorname{dist}(Z,\overline{P_x})<\eta\}.
\end{equation}
Consequently, if $t_n\downarrow0$ and paths in $D_{t_n}^x$ start at
$p$, stay in one fixed bounded box, and have endpoints tending to $z$,
then $z\in\overline{P_x}$.  The same holds for starting points tending
to $p$ inside a uniform initial ball.  Fixed auxiliary parameters are
permitted; generic restrictions are made in the real centers used for
the recenterings.
\end{theorem}

The proof uses the recentering identity in \eqref{loc:normalized}
to turn paths at nearby real centers into real families of points
in one normalized domain.  In the limit these families have the
form $u+G_x(u)F(a)$, where $F$ is a limiting path.  We will attach
analytic discs to them, solving from the last block to the first,
and obtain centers in the marked kernel approaching the path
endpoint.  The next two lemmas supply replacement paths whose
dependence on the real center and on the path parameter permits
this disc construction.

\subsection{Controlled paths from ordinary uniformization}

We use two precise forms of real-analytic resolution: a closed
subanalytic set has a proper real-analytic uniformization by a manifold
of the same dimension, and a nonzero analytic function becomes a
monomial times a nonvanishing analytic function after a proper
surjective real-analytic map.  These are Theorem~0.1 and Corollary~4.9 of
\cite{BierstoneMilman1988}.  The relative estimates below are consequences
proved here, rather than an additional relative resolution theorem.

\begin{lemma}[Generic analytic power parametrization]
\label{loc:puiseux}
Let $U\subset\mathbb R^k$ be open and let
$f:U\times(0,\epsilon)\to\mathbb R^m$ be bounded and globally subanalytic.
Outside a closed subanalytic subset of $U$ of dimension less than $k$,
locally in $x$, there are an integer $q\geq1$ and a real-analytic map
$a(x,s)$, defined near $U'\times\{0\}$, such that
\[
                 f(x,t)=a(x,t^{1/q})\qquad(0<t<\epsilon').
\]
Here $U'$ and $\epsilon'$ may be reduced.  If a scalar coordinate of $f$
is positive, it either has a nonzero limit or, after a further generic
restriction, has the form
\[
                 t^{\nu/q}a_0(x,t^{1/q}),\qquad a_0(x,0)>0,
                 \quad \nu\in\mathbb N.
\]
On smaller parameter boxes the analytic maps have a common complex
neighborhood and all their derivatives are bounded.
\end{lemma}

\begin{proof}
Apply Lemma~\ref{prep:puiseux} to the finitely many coordinate
functions of $f$.  On a generic parameter box they have joint
convergent Puiseux expansions with a common denominator.  Boundedness
forces every negative-power coefficient to vanish, giving the
analytic map $a(x,t^{1/q})$.  For a positive coordinate, restrict
further so that its first nonzero coefficient does not vanish.
Positivity then makes that coefficient positive and gives the
asserted factorization.  The discarded preparation cells and
coefficient zero loci come from finitely many data on the bounded
boxes in use.  Their lower-dimensional relative closures give the
closed subanalytic exceptional set in the statement.  Restricting
the joint analytic extensions to smaller compact parameter boxes
gives a common complex neighborhood and bounds for every derivative.
\end{proof}

\begin{lemma}[Uniformly controlled replacement paths]
\label{loc:paths}
Let $U\subset\mathbb R^k$ be open, let
$D\subset U\times(0,\epsilon)\times\mathbb R^n$ be a bounded
globally subanalytic family of open sets, and let $v,w$ be bounded
subanalytic sections that belong to the same component of $D_{x,t}$.
Locally at generic $x_0\in U$, there are a neighborhood $U'$ of $x_0$,
$\epsilon'>0$, and paths
\[
 F_t:U'\times[0,1]\longrightarrow\mathbb R^n
\]
with endpoints $v(x,t),w(x,t)$ and image in $D_{x,t}$, with the following
uniform properties.  For each $t$, $F_t(\cdot,a)$ extends holomorphically
to one fixed complex neighborhood $U'_{\mathbb C}$, independently of
$t$ and $a$.  On a smaller such neighborhood there is $C<\infty$ such
that
\begin{equation}\label{loc:path-estimate}
 \sup_{0<t<\epsilon'}\sup_{x\in U'_{\mathbb C}}
 \bigl(|F_t(x,a)|+|F_t(x,a)-F_t(x,b)|/|a-b|\bigr)\leq C
 \quad(a\ne b).
\end{equation}
In particular, every fixed order of $x$-derivatives satisfies the same
estimate, with its own constant.  The replacement paths and their
timing need not be subanalytic.
\end{lemma}

\begin{proof}
Apply definable Hardt triviality to the projection of $D$ onto $(x,t)$,
compatibly with the two marked endpoint graphs
\cite[Theorem~5.22]{Coste1999}. On each base cell
the trivialization identifies these sections with fixed points in a
model fiber.  They lie in the same component; choose a fixed definable
path between them and transport it by the continuous trivialization.
This gives continuity jointly in $(x,t,a)$ on that cell.  The finite
base walls have dimension at most $k$, and their limiting loci at
$t=0$ have dimension less than $k$.  Deleting those loci puts a smaller
product $U\times(0,\epsilon)$ in one cell.  Retain the original
path parameter $a$ and let $T$ be the closure of its graph in
$(x,t,a,z)$.  This graph has dimension $k+2$.  Over $t>0$ and interior
base parameters its closure is still that graph.  Uniformize $T$
properly, resolve $t$, and discard components supported at $t=0$.
The maps from the resulting manifold into $(x,t,a,z)$ are analytic.

For every stratum of the normal-crossings divisor, delete the closure
of its critical values in $x$.  A stratum that cannot dominate the
$k$-dimensional base is avoided altogether.  At every remaining point
over $(x_0,0)$, the original $x$ coordinates can be retained as
tangential coordinates while all divisor coordinates are retained as
normal coordinates.  Indeed the differential of $x$ is surjective
on their common zero stratum, so the analytic inverse function theorem
allows precisely this coordinate change.  Absorbing the nonzero
monomial unit gives, on each fixed sign piece,
\begin{equation}\label{loc:monomial}
                t=y_1^{m_1}\cdots y_r^{m_r},\qquad y_i>0.
\end{equation}
The coordinates not displayed are free coordinates in an ordinary
box.  We use smaller closed boxes inside larger analytic chart boxes.
Properness gives finitely many of them that cover every relevant
graph point for $x$ near $x_0$ and $t$ small.  Their maps have common
complex neighborhoods in $x$ after further shrinking.

Here is a precise gluing argument.  At fixed $(x,t)$, the image of a
smaller closed chart box with fixed signs and
\eqref{loc:monomial} is compact and connected: in logarithmic
coordinates its monomial part is an affine hyperplane intersected
with coordinate half-spaces, and the free part is a box.
The positivity constraints do not spoil compactness: if $y_h\leq B_h$
and $t>0$ is fixed, then
$y_i\geq(t/\prod_{h\ne i}B_h^{m_h})^{1/m_i}>0$.
These finitely many sets are subsets of the actual continuous path
graph for $t>0$ and cover that graph.  The
intersection graph of the sets needed to connect its two endpoints
therefore has a chain joining them; use a simple chain, whose length
is bounded by the number of chart/sign pieces.  The existence of
each finite chain, its intersection points, and their lifts is a
subanalytic condition.  Selection followed by a generic restriction
fixes one chain and bounded subanalytic endpoint lifts for all
$(x,t)$ in a smaller product.  Adjacent endpoint images agree exactly.
All endpoint lifts remain in the larger analytic boxes.  This argument
uses closed boxes to retain the interface points and does not assume
that open chart images overlap uniformly.

By Lemma~\ref{loc:puiseux}, all these finitely many endpoint
coordinates have analytic power parametrizations with a common
denominator.  For two positive endpoint coordinates $A_i(x,t)$ and
$B_i(x,t)$, join them by
\[
 y_i(x,t,s)=A_i(x,t)^{1-s}B_i(x,t)^s,
                   \qquad 0\leq s\leq1.
\]
The monomial equality is retained, since it holds at both endpoints.
The interpolation stays in the coordinate box.  Free coordinates are
interpolated linearly.  If needed, insert one interior waypoint before
making these interpolations; this also retains specified strict sign
conditions.  Every image point with $t>0$ still belongs to the original
graph, and hence its $z$ coordinate belongs to $D_{x,t}$.

For completeness, the clock controlling the speed is independent of
$x$.  Write $y_i^0(s)=y_i(x_0,t,s)$.  Positive endpoint coordinates have
the form $t^\alpha$ times analytic positive units.  On a fixed small
complex neighborhood of $x_0$, logarithms of the unit ratios to their
values at $x_0$ are analytic and uniformly bounded, with bounded
derivatives.  It follows that
\[
 y_i(x,t,s)=y_i^0(s)E_i(x,t,s),\qquad
 |E_i|+|\partial_s E_i|\leq C,
\]
uniformly also for complex $x$.  Thus
\[
 |\partial_s y_i(x,t,s)|
           \leq C\bigl(|(y_i^0)'(s)|+y_i^0(s)\bigr).
\]
Each $y_i^0$ is monotone and bounded.  Define the increasing clock
\[
 h_t(s)=
 \frac{s+\sum_i|y_i^0(s)-y_i^0(0)|}
 {1+\sum_i|y_i^0(1)-y_i^0(0)|}.
\]
After $s=h_t^{-1}(a)$, the displayed derivative estimate, the bounded
total variations of the $y_i^0$, and the elementary estimates for the
free coordinates give a common Lipschitz bound.  Shrinking the complex
$x$ neighborhood ensures that the complex interpolations remain in the
complex analytic chart neighborhoods.  Cauchy estimates give the same
Lipschitz bounds for every $x$-derivative on a smaller neighborhood.
Analytic chart maps preserve the estimates.  Allocate fixed consecutive
subintervals to the finitely many chart segments and concatenate them.
The joints agree for real $x$, and hence also on the complex
neighborhood by analytic uniqueness.  This proves
\eqref{loc:path-estimate}.
\end{proof}

\subsection{Attached discs and upper containment}

\begin{proof}[Proof of Theorem~\ref{loc:theorem}]
We first prove the analytic assertion for one controlled path family
provided by Lemma~\ref{loc:paths}.  Fix its generic real center $x$ and
suppress $x$ on $R_t,G$ and $P$.  Suppose
$F_t(x,0)=p$ and the paths are uniformly bounded.  Passing to a
subsequence, the uniform Lipschitz estimate gives
\[
                  F_t(x,\cdot)\longrightarrow F(\cdot)
                        \quad\hbox{in }C^\alpha([0,1])
                        \quad(0<\alpha<1).
\]
Here and below this convergence is along the selected subsequence.
The derivative bounds in Lemma~\ref{loc:paths} and $R_t\to0$ show that
on every fixed finite real $u$ box the actual graphs
\begin{equation}\label{loc:actual-graphs}
 \Phi_t(u,a)=u+C_t(x,u)F_t(x+R_t(x)u,a)
\end{equation}
converge to
\begin{equation}\label{loc:limit-graphs}
                       \Phi(u,a)=u+G(u)F(a).
\end{equation}
This holds in $C^\alpha$ in $a$, with any fixed finite number of
$u$ derivatives, after lowering $\alpha$ if necessary.  Every real
point of \eqref{loc:actual-graphs} belongs to $D_t^x$: in original
coordinates it is exactly
\[
 x'=x+R_t(x)u,\qquad x'+R_t(x')F_t(x',a)\in S.
\]
This identity is why all bounded real displacements, rather than a
small fixed collection of orbits, must be retained.

The initial graph has room at every real displacement needed below.
Indeed, the initial kernel condition, definable selection, and generic
restriction give a real neighborhood $U'$ and numbers $r,\epsilon>0$
such that
\begin{equation}\label{loc:uniform-start}
 B(p,r)\subset D_t^{x'}
       \qquad(x'\in U',\ 0<t<\epsilon).
\end{equation}
To see the uniformity, select positive radii and thresholds definably
from the initial kernel condition, restrict to a cell on which they
are continuous, and shrink to a relatively compact smaller box.
For every fixed $M$, \eqref{loc:ratio} and invertibility give $c_M>0$
such that
\begin{equation}\label{loc:orbit-room}
 B\bigl(u+C_t(x,u)p,c_Mr\bigr)\subset D_t^x
               \qquad(|u|\leq M,\ 0<t<\epsilon_M).
\end{equation}
One may take $c_M$ smaller than
$\bigl(2\sup_{|u|\leq M}\|G(u)^{-1}\|\bigr)^{-1}$.
Thus the later real corrections are admissible even when they do not
stay close to zero.

\emph{The triangular equation.}
We use Bishop's reduction of the attached-disc condition to a
boundary conjugation equation \cite{Bishop1965}; see also
\cite[\S2, equation~(2.1) and Proposition~2.1]{HillTaiani1978}.
The triangular polynomial system and the endpoint estimates needed
here are established below. Normalize circle measure to have total mass one, denote its mean by
$\langle\cdot\rangle$, and let $H$ be the real Hilbert transform with
$H(1)=0$ and $H(\cos n\theta)=\sin n\theta$.  A complex boundary value
$v+iw$ bounds a holomorphic disc precisely when
$v-\langle v\rangle=-Hw$.  The operator $H$ is bounded on $C^\alpha$
for $0<\alpha<1$ by the conjugate-function estimate of
\cite[pp.~100--102]{Priwaloff1916}, and on $L^q$ for $1<q<\infty$
by the M.~Riesz theorem in the form proved in
\cite[equation~(0.1) and \S\S1--2]{Essen1984}.

Choose a Lipschitz schedule $a:\partial\mathbb D\to[0,1]$ that equals
zero on a nonempty open arc $I$.  Choose a smooth nonnegative function
$\chi$ supported in $I$, with $\langle\chi\rangle>0$.  Add a complex
parameter $\zeta_j\chi$ in each block.  Write
\[
 V_j(u,\theta,\zeta)
     =G_j(u_{>j}(\theta))F_j(a(\theta))+\chi(\theta)\zeta_j.
\]
The real unknowns $u_j\in C^\alpha(\partial\mathbb D,\mathbb R^{d_j})$
are required to have mean zero.  Their equations are exactly
\begin{equation}\label{loc:bishop}
 u_j=-H\operatorname{Im}V_j
       -(I-\langle\cdot\rangle)\operatorname{Re}V_j,
                       \qquad j=s,s-1,\ldots,1.
\end{equation}
Starting with block $s$, whose matrix is the identity, these formulas
define a unique solution successively.  They give a holomorphic disc
$f_{a,\zeta}$ with boundary $u+G(u)F(a)+\chi\zeta$.  Polynomial
composition is continuous on $C^\alpha$, so every compact family of
schedules and sufficiently small perturbations has a uniform
$C^\alpha$ bound on its solutions.  In particular it uses just one
finite real $u$ box.  No estimate is asserted uniformly over schedules
whose transition intervals subsequently shrink to zero.

Its center is
\begin{equation}\label{loc:center}
 C_j(a,\zeta)=\langle G_j(u_{>j})F_j(a)\rangle
                       +\langle\chi\rangle\zeta_j.
\end{equation}
Later blocks do not depend on earlier perturbations.  Hence the real
differential of $\zeta\mapsto C(a,\zeta)$ is block triangular, with
diagonal $\langle\chi\rangle I$ on each complex block.  It is invertible
at every solution.  This proves full real rank $2d$, not only motion
in the imaginary center directions.

\emph{Persistence and filling.}
Replace $G(u)F(a)$ in the equations by
$C_t(x,u)F_t(x+R_t(x)u,a)$.  On the Banach space of mean-zero real
$C^\alpha$ functions, the differential of the limiting left-hand
system is identity plus a strictly block triangular bounded operator.
Its inverse is the finite sum for a nilpotent triangular operator.
The convergence after \eqref{loc:actual-graphs} gives convergence of
these systems and their first derivatives in operator norm near any
compact family of limiting solutions.  For example it follows
directly from the $C^2$ bounds in $u$, the product estimate in
$C^\alpha$, and convergence of the $a$-dependent coefficients in a
slightly stronger H\"older norm.  The implicit function theorem
therefore gives nearby solutions, uniformly along the compact
family, and convergence of their centers and center differentials.
Local uniqueness patches the solutions over the family.  To justify a
uniform center-image ball, let $K$ bound the inverse differentials of
the limiting center maps at zero perturbation.  Compactness and $C^1$
convergence allow a fixed perturbation radius $r_0>0$ on which the
variation of those differentials, including the small-$t$ error, is
at most $1/(4K)$, while the small-$t$ inverse differentials at zero
have norm at most $2K$.  The contraction proof of the inverse function
theorem then puts a ball of radius $r_0/(4K)$ about each unperturbed
center in its center-map image.  One first decreases $r_0$ further if
the start-room restriction below requires it.

Use the homotopy of schedules $a_\sigma=\sigma a$, $0\leq\sigma\leq1$;
all retain the start arc $I$.  First use a preliminary closed
perturbation ball to bound the limiting solutions in a finite $u$ box.
Enlarge that box slightly, compute its start-room constant in
\eqref{loc:orbit-room}, and then decrease the perturbation ball so that
$\|\chi\zeta\|_\infty$ is smaller than half that room.  Only after these
choices take $t$ small.
On the support of $\chi$, \eqref{loc:orbit-room} admits all the chosen
small complex perturbations.  Away from that support the boundary
lies on the actual graphs \eqref{loc:actual-graphs}.  Thus the entire
compact homotopy of disc boundaries lies in $D_t^x$.  At $\sigma=0$
and $\zeta=0$ the limiting solution is $u=0$ and the disc is constant
$p$.  For small $t$ and small $\zeta$, its actual disc lies in the
initial ball.  The continuity principle in the locally pseudoconvex
domain now keeps every filled disc in $D_t^x$.

There is no issue at the exterior of the original coordinate box:
for this fixed schedule the normalized boundary values, and hence
the entire discs, are bounded by a fixed constant.  Their original
coordinates converge uniformly to $x\in\Omega$.  Choose a fixed
convex coordinate box $\Omega_0$ with $x\in\Omega_0\Subset\Omega$.
The intersection $S\cap\Omega_0$ is locally pseudoconvex at every
boundary point in $\mathbb C^d$, using the given condition on $S$ and
convexity at the artificial boundary.  Its components are therefore
Stein by the Euclidean Levi theorem.  A plurisubharmonic exhaustion
of the component in use bounds every filled disc by its boundary
maximum; the compact boundary homotopy prevents a first exit.  All
the discs stay in $\Omega_0$ by their already established coordinate
bounds.

For each $\sigma$, the uniform center rank supplies a fixed ball
$B(C(a_\sigma,0),\rho)$ contained in $D_t^x$ along the chosen
subsequence, after slightly decreasing $\rho$.  This already implies
eventual inclusion for all small $t$: the set
\[
 \{t>0:B(C(a_\sigma,0),\rho)\subset D_t^x\}
\]
is a one-dimensional subanalytic set and has $0$ in its closure.
It therefore contains an interval $(0,\epsilon_\sigma)$.
Consequently $C(a_\sigma,0)$ belongs to the full open kernel.  The
continuous curve of these centers starts at $p$, so all its points
belong to $P$.  This step does not require the replacement paths or
the schedule homotopy themselves to be subanalytic.

\emph{Approaching the endpoint.}
For each $\delta>0$, choose a Lipschitz loop $a_\delta$ that starts at
$0$, goes to $1$ and returns, equals $0$ on a nonempty arc, and equals
$1$ outside a set of measure at most $\delta$.  Apply the preceding
argument at this fixed $\delta$, with zero perturbation.  It gives
$C(a_\delta,0)\in P$.  Set $c=F(1)$.  Since $F$ is bounded,
\[
 F(a_\delta)\longrightarrow c\quad\hbox{in }L^q
                            \quad(1<q<\infty).
\]
We claim that every block of $u^\delta$ from \eqref{loc:bishop}
converges to zero in every such $L^q$.  For the last block this is
the boundedness of $H$ and subtraction of the mean.  If it holds
for the later blocks, polynomiality gives
$G_j(u_{>j}^\delta)\to I$ in every finite $L^q$: each monomial is
estimated by H\"older's inequality using the finitely larger
exponents required by its degree.  Multiplication by the uniformly
bounded $F_j(a_\delta)$ then gives $V_j\to c_j$ in every finite
$L^q$.  Equation~\eqref{loc:bishop} proves the induction.  Taking
means in \eqref{loc:center} yields
\[
                         C(a_\delta,0)\longrightarrow F(1).
\]
Thus $F(1)\in\overline P$.  The order of choices was: fix the loop
and its nonempty start arc; solve and bound all corrections; choose
their finite real box; let $t\downarrow0$ and obtain a center in $P$;
only then shrink the exceptional arcs.  No shrinking-arc estimate is
used to justify the finite-$t$ implicit function theorem.

\emph{The generic and uniform quantifiers.}
It remains to pass from controlled paths to all paths.  The family
of kernels $\mathcal K(D_t^x)$ is subanalytic, by its quantified ball
definition; so is its component $P_x$ marked by $p$.  For fixed integer
$L>|p|$ and positive rational $\eta$, let $E_{L,\eta}$ be the set of
real centers for which arbitrarily small $t$ admit a point of
$A_t^x(L)$ at distance at least $\eta$ from $\overline{P_x}$.
This is a subanalytic set, using definability of components.

If $E_{L,\eta}$ had interior, choose within it a generic smaller real
box.  One-dimensional definability in $t$ and selection give endpoints
and paths to them for every sufficiently small $t$, after reducing
the box generically.  Use a slightly larger fixed spatial box for the
paths.  Lemma~\ref{loc:paths} supplies a controlled replacement family
there.  At a generic fixed center, a subsequence of its endpoints
converges, and the analytic argument just proved puts its limit in
$\overline{P_x}$.  This contradicts their distance of at least $\eta$.
Thus $\dim E_{L,\eta}<d$.  Delete these exceptional sets for the
countably many $L,\eta$.  At every remaining center the negation of
membership in $E_{L,\eta}$ is exactly \eqref{loc:upper}.  Increasing
integer boxes and decreasing rational accuracies give the stated
quantifiers for all finite boxes and all positive accuracies.

Finally, a starting point approaching $p$ can be joined to $p$ by a
segment in the uniform initial ball.  Enlarging the fixed path box
by a fixed amount accommodates this segment and proves the last
assertion.
\end{proof}

\begin{corollary}[Semidefinite outer tests]\label{loc:matrix-outer}
Suppose, in addition, that the marked kernel is described recursively
by strict positive-definiteness tests for finitely many continuous
Hermitian matrix functions $M_j(Z)$, and that every point of
$\overline P$ satisfies $M_j(Z)\geq0$.  Then for each finite $L$,
at a generic fixed center,
\[
 \inf_{Z\in A_t(L)}\lambda_{\min}(M_j(Z))\geq-o(1)
                \qquad(t\downarrow0)
\]
simultaneously for all $j$.  The assertion is uniform over the
accessible component in that fixed box.
\end{corollary}

\begin{proof}
If one inequality failed, a sequence in the bounded box would have
a convergent subsequence on which the corresponding least eigenvalue
is bounded above by a fixed negative number.  Theorem~\ref{loc:theorem}
puts its limit in $\overline P$, contradicting continuity and the
semidefinite test there.  There are only finitely many tests.
\end{proof}

The corollary is applied first at a fixed radius.  A diagonal choice
of radii and scales may subsequently impose these outer estimates
together with compact inner inclusion.  In particular, an arbitrary
joint limit of the radius and scale is not part of the assertion.
